\documentclass[12pt,a4paper,oneside]{article}
\usepackage[brazil]{babel}
\usepackage[utf8]{inputenc}
\usepackage{olymp}
\usepackage{multirow}
\usepackage[table]{xcolor}

\setcounter{problem}{0}

\begin{document}

\begin{problem}{Vendas do funcionário do mês}{vendas}{2}

Uma empresa tem uma tabela com o relatório das vendas de cada funcionário em cada dia do mês.
Por exemplo, a seguinte tabela mostra as vendas dos $3$ funcionários em $4$ dias.

\begin{center}
\begin{tabular}{cc|ccc|}
    & & \multicolumn{3}{c|}{Funcionário} \\ %\cline{3-5}
    & & 1 & 2 & 3 \\ \hline
    \multirow{4}{*}{\rotatebox{90}{Dia}}
    & 1 & \cellcolor{gray!20}10 & \cellcolor{gray!20}5  & \cellcolor{gray!20}7  \\
    & 2 & \cellcolor{gray!20}20 & \cellcolor{gray!20}8  & \cellcolor{gray!20}11 \\
    & 3 & \cellcolor{gray!20}14 & \cellcolor{gray!20}7  & \cellcolor{gray!20}12 \\
    & 4 & \cellcolor{gray!20}13 & \cellcolor{gray!20}18 & \cellcolor{gray!20}15 \\ \hline
    \multicolumn{2}{c|}{Total} & 57 & 38 & 45
\end{tabular}
\end{center}

A empresa planeja premiar o ``funcionário do mês'', que é o que efetuou maior total de vendas no mês.
Como primeiro passo, ela quer saber apenas as vendas do funcionário do mês. No exemplo acima, $57$.

%\Notes
%
%Observações, se tiver.

\InputFile

A primeira linha contém dois inteiros, $D$ e $F$, que são o número de dias e de funcionários que constam no relatório.
Em seguida existem $D$ linhas, uma para cada dia. Cada uma delas contém $F$ inteiros, indicando o total de vendas de cada funcionário naquele dia.

Restrições:
$2 \leq D \leq 30$
e $2 \leq F \leq 20$.

\OutputFile

Seu programa deve informar a maior venda total feita por um funcionário.

%\Notes

\Example

\begin{example}
\exmp{
4 3
10 5 7
20 8 11
14 7 12
13 18 15
}{
57
}%
\end{example}

\begin{example}
\exmp{
3 3
10 20 30
40 50 60
70 80 90
}{
180
}%
\end{example}

\begin{example}
\exmp{
2 4
7 10 4 5
5 4 8 9
}{
14
}%
\end{example}
%Comentários sobre o exemplo, se necessário.

\end{problem}

\end{document}